\section{Origin of matter: topologically locked phase defects and inertial budget}
\label{sec:matter_defects}

\subsection{Topological charge from curvature flux (Dirac/Wu--Yang)}
\label{subsec:topo_charge}

CAP treats stable localized ``matter'' as a defect sector of the compensating connection introduced in Section~\ref{sec:gauge_fields}. The relevant rigidity statement is standard in gauge theory and topology: nontrivial defect sectors are characterized by integral cohomology classes (Chern numbers) extracted from curvature flux \cite{Dirac1931,WuYang1975,Nakahara2003,Mermin1979}.

\paragraph{U(1) example.}
Let $U=\RR^{3}\setminus\{0\}$ be space with an isolated puncture. A $U(1)$ gauge field is described by local 1-forms $A$ with curvature 2-form $F=\dd A$, which satisfies the Bianchi identity $\dd F=0$. Surround the defect by a sphere $S^{2}\subset U$. Global consistency of phase (single-valued transition functions between local gauges) implies the first Chern number is an integer:
\begin{equation}
  Q_{\mathrm{top}}
  :=\frac{1}{2\pi}\int_{S^{2}} F
  \in \ZZ.
\end{equation}

\begin{lemma}[``Closed but not exact'' on a punctured region]
\label{lem:closed_not_exact}
On $U=\RR^{3}\setminus\{0\}$, a smooth 2-form $F$ with $\dd F=0$ and $\int_{S^{2}}F=2\pi n\neq 0$ cannot be exact. Equivalently, no global gauge choice can remove the defect sector.
\end{lemma}
\begin{proof}
If $F=\dd A$ globally on $U$, then Stokes' theorem gives $\int_{S^{2}}F=\int_{S^{2}}\dd A=0$, contradicting $\int_{S^{2}}F=2\pi n\neq 0$. In de Rham language, the flux labels the nontrivial class in $H^{2}(U)\cong H^{2}(S^{2})\cong \ZZ$ \cite{Nakahara2003}.
\end{proof}

\subsection{Mass as sustained budget: internal winding and geometric impedance}
\label{subsec:mass_budget}

Omega Dynamics further proposes an organizing hypothesis: stable matter corresponds to phase errors that cannot be removed by local gauge redefinitions (closed but not exact), i.e.\ \emph{topologically locked defects}. Maintaining a nontrivial topological sector requires persistent implementation budget, which manifests macroscopically as inertial mass and as resistance to propagation (geometric impedance).

In HPA language, such defects can be modeled as localized impedance centers in the scan--phase dynamics: routing around a local algebraic obstruction requires extra scan cycles and accumulates additional geometric phase. From the readout perspective, this appears as propagation delay and curvature, consistent with the computational-lapse dictionary of Section~\ref{subsec:two_times}.

\paragraph{A quantitative interface inequality.}
The phase-potential closure (Section~\ref{subsec:phase_potential}) assigns to a nonnegative coarse-grained mismatch cost density $\sigma$ a monopole strength
\begin{equation}
  M := \kappa_\Phi \int_{\RR^3}\sigma(\mathbf{x})\,\dd^3\mathbf{x},
\end{equation}
so that $\Phi(r)\sim -M/r$ for isolated sources. A topologically nontrivial gauge sector (Lemma~\ref{lem:closed_not_exact}) necessarily carries nonzero curvature and therefore incurs a positive maintenance cost in any local quadratic effective action. In soliton models this appears as Bogomol'nyi-type bounds that lower-bound energy by topological charge \cite{MantonSutcliffe2004,Rajaraman1982}. CAP packages the same structural fact as a minimal, testable interface inequality:
\begin{proposition}[Topological lower bound on mismatch charge]
\label{prop:topo_mass_bound}
Assume sustaining a defect sector with $|Q_{\mathrm{top}}|$ requires mismatch charge $Q_\sigma:=\int_{\RR^3}\sigma\,\dd^3x$ satisfying
\begin{equation}
  Q_\sigma \ge m_{0}\,|Q_{\mathrm{top}}|
\end{equation}
for some microscopic per-charge budget $m_0>0$. Then the exterior monopole coefficient obeys
\begin{equation}
  M \ge \kappa_\Phi\, m_{0}\,|Q_{\mathrm{top}}|.
\end{equation}
\end{proposition}

At the current stage, $m_0$ and the full defect spectrum (charges, stability, interactions) must be computed from explicit microscopic scan/readout architectures; this remains an open task within the broader $\Omega$ program \cite{Ma2025FoP}.


